% Canberk Sönmez

\documentclass{article}
\usepackage[margin=2cm]{geometry}
\usepackage{amsmath}
\usepackage{graphicx}
\usepackage{lipsum}
\usepackage{hyperref}
\usepackage{discussionbox}
\renewcommand{\familydefault}{\sfdefault}

\DefinePerson{Intro}{black}
\DefinePerson{Basics}{black}
\DefinePerson{Communication}{black}

\begin{document}

\begin{discussionbox}[Tutorial (D\dbox{nextid})]
    \Intro{
        DiscussionBox provides a forum-like typesetting for discussions taking place on \LaTeX{} document sharing platforms like Overleaf.
        Compared to review mechanisms provided by such platforms, DiscussionBox enables non-volatile and structured conversations, somewhat to e-mail.
        Unlike e-mail, however, you can benefit from powerful \LaTeX{} features!
    }

    \begin{dboxmsg}{Basics}
        The fundamental component of the DiscussionBox package is, well, a \texttt{discussionbox}:
        \begin{verbatim}
    % discussion boxes are declared as an environment:
    \begin{discussionbox}[Discussion \dbox{nextid}] % title is optional
        % you can use \dbox{nextid} to enumerate titles
        % the default title is "Discussion Box #1"

        % messages...
    \end{discussionbox}
        \end{verbatim}
    \end{dboxmsg}

    \begin{dboxmsg}{Basics}
        A discussion box is merely a container for \textit{message}s and a message is typed by a \textit{person}.

        \begin{verbatim}
    % in the preamble, define a person:
    \DefinePerson{Canberk}{blue} % \DefinePerson{name}{color}
    
    % within document:

    \begin{discussionbox}

        % messages using a single-name command:
        \Canberk{
            Your message here...
        }

        % or, in an environment: (which might work better in some cases)
        \begin{dboxmsg}{Canberk} % \begin{dboxmsg}{person name}...\end{dboxmsg}
            Your message here...
        \end{dboxmsg}

    \end{discussionbox}
        \end{verbatim}
    
    \noindent The name of the author is shown in the beginning of the message.
    Furthermore, the entire message colored with the color of the author.\\
    Within a message, you are allowed to use all the \LaTeX{} features, you can look for the examples below.
    \end{dboxmsg}

    \Communication{\LabelMsg{msg:msg1}
        The main idea of a DiscussionBox is to enable effective communication over.
        For that purpose, it provides three important features: (1) \textit{message addresses}, (2) \textit{references} and (3) \textit{replies}.
    }
    \begin{dboxmsg}{Communication}
        Each message in a discussion box is uniquely identified by an address.
        For example, the address of this one \dbox{msgaddr}.
        D$x$ stands for the discussion with number $x$ and the rest of the address identifies the message within a discussion.
        Addresses appear right next to the author name and formatted as superscript.
        You can use these addresses to refer to messages.
        Similar to other \LaTeX{} objects that can be referenced, you need to first label the message:

        \begin{verbatim}
    % use \LabelMsg{label} to label a message
    \Canberk{Hello, I say things .... \LabelMsg{canberk:things}}

    % later...

    \Mehmet{Note what Canberk says \ref{canberk:things}, he says some serious stuff.}
        \end{verbatim}
        
        \noindent Let's reference the previous message: \ref{msg:msg1}!
        You can click on the reference to navigate to the message.
        \LabelMsg{msg:msg2}

        \begin{dboxmsg}{Communication}
            A message can be a reply to another message.
            For example, this message (\dbox{msgaddr}) replies to \ref{msg:msg2}.
            Message addresses reflect reply-to relations.
            If a message has address D$x$.$y$.$z$.$n$, it is the $n$\textsuperscript{th} reply to D$x$.$y$.$z$.
        \end{dboxmsg}
        \begin{dboxmsg}{Communication}
            For instance, the last character of the address of this message is 2.
            A reply message is nested within the message that it replies to:
            \begin{verbatim}
    \Mehmet{I think Switzerland is the most prominent economy in Europe.
            \Canberk{I disagree, it is Germany.}
    }

    % works the same way with environments
            \end{verbatim}
        \end{dboxmsg}
    \end{dboxmsg}
\end{discussionbox}

\DefinePerson{Canberk}{magenta}
\DefinePerson{Ahmet}{blue}
\DefinePerson{Mehmet}{red}

\begin{discussionbox}
    \Canberk{
        Hey, what's up?
        \Ahmet{
            It's perfect! By the way, math formatting works:
            \begin{align}
                F = ma
            \end{align}
            \Ahmet{Wow, that's cool!}
        }
        \Mehmet{
            Same here! Check out the following Latex logo:

            \begin{center}
                \includegraphics[width=0.3\linewidth]{fig/latex-logo.png}
            \end{center}
            \Ahmet{
                Wow, that's so good!
            }
            \Canberk{
                I totally agree!
                \Canberk{\LabelMsg{msg:canberk} More and more nesting...}
            }
            Previous message continues...
            \Ahmet{Here is a long text:\\ \lipsum[1-2]}
        }
    }
\end{discussionbox}

\begin{discussionbox}
    \Ahmet{
        \lipsum[1]
        \Canberk{
            \lipsum[2-3]
            \Mehmet{
                \lipsum[4-5][3]
            }
        }
    }
    \Ahmet{
        \lipsum[6-7]
        \Canberk{
            \lipsum[8-9]
            \Mehmet{
                Here is the reference \ref{msg:canberk}.
                \lipsum[10-11]
            }
        }
    }
\end{discussionbox}


\end{document}
